
\documentclass{article}
\usepackage{colortbl}
\usepackage{makecell}
\usepackage{multirow}
\usepackage{supertabular}

\begin{document}

\newcounter{utterance}

\centering \large Interaction Transcript for game `dond', experiment `semi\_en', episode 19 with reflexion{-}Qwen3.5{-}4B.
\vspace{24pt}

{ \footnotesize  \setcounter{utterance}{1}
\setlength{\tabcolsep}{0pt}
\begin{supertabular}{c@{$\;$}|p{.15\linewidth}@{}p{.15\linewidth}p{.15\linewidth}p{.15\linewidth}p{.15\linewidth}p{.15\linewidth}}
    \# & $\;$A & \multicolumn{4}{c}{Game Master} & $\;\:$B\\
    \hline

    \theutterance \stepcounter{utterance}  
    & & \multicolumn{4}{p{0.6\linewidth}}{
        \cellcolor[rgb]{0.9,0.9,0.9}{
            \makecell[{{p{\linewidth}}}]{
                \texttt{\tiny{[P1$\langle$GM]}}
                \texttt{You have attempted this same task before and did not succeed. These are the plans you wrote after those attempts.} \\
\texttt{Trial \#1: New plan: The last attempt was aborted because I included explanatory text and a question ("Do you agree...?") in my first message, violating the rule that non{-}proposal messages must be direct transmissions without extra commentary or requests for agreement. I must strictly separate negotiation from the formal proposal. My strategy should be: First, send a concise statement of my valuation and a clear, direct offer of a specific partition (e.g., "I take the statue and one whistle; you take the comb and the other whistle") without asking for their opinion. If they accept verbally, we proceed to the secret proposal phase. If not, I will adjust my offer in subsequent messages to be more attractive (e.g., offering the statue and both whistles) while maintaining the strict format of direct statements only. Finally, when ready to finalize, I will use the exact secret proposal format `{[}Proposal: <number> <object name>, ...{]}` in a single message without any surrounding text.} \\
\\ 
\texttt{Trial \#2: New plan: The previous attempt failed because I made a secret proposal (`{[}Proposal: 1 statue, 2 whistles{]}`) that was not complementary to the opponent's likely counter{-}proposal. Since the opponent also values the statue highly (8) and whistles moderately (1), they would likely claim the statue and one or both whistles, leaving me with only the comb (value 0) and the wrench (value 0), resulting in a total score of 0. My mistake was proposing a bundle that the opponent could not reasonably accept without giving up their high{-}value item. The correct strategy is to first negotiate verbally to reach a mutually acceptable split (e.g., I take the statue and one whistle; you take the comb and the other whistle, or I take the statue and the wrench; you take the comb and both whistles) to ensure both parties get a positive score. Only after confirming a verbal agreement should I submit a secret proposal that exactly matches this agreed{-}upon split, ensuring the proposals are complementary and both players receive a non{-}zero score.} \\
\\ 
\texttt{Trial \#3: New plan: The last attempt failed because I submitted a secret proposal (`{[}Proposal: 1 statue, 1 whistle{]}`) before establishing a verbal agreement with the opponent. According to the rules, proposals are only valid if they are complementary to the other player's proposal, and since I had not yet received a counter{-}proposal, the system could not verify complementarity, leading to an abort. My strategy must be: First, send a direct verbal offer of a specific item split (e.g., "I take the statue and one whistle; you take the comb and the other whistle") without using square brackets or any explanatory text. Wait for the opponent to verbally confirm or counter this offer. Only after receiving a clear verbal agreement (or a counter{-}offer that I can accept) should I submit the final secret proposal using the exact format `{[}Proposal: <number> <object name>, ...{]}` to lock in the agreed{-}upon division.} \\
\\ 
\\ 
\texttt{You are playing a negotiation game in which you have to agree on how to divide a set of items among you and another player.} \\
\\ 
\texttt{Rules:} \\
\texttt{(a) You and the other player are given a set of items. Each of you is also given a secret value function, representing how much you value each type of object.} \\
\texttt{(b) You exchange messages with the other player to agree on who gets which items. You can send a maximum of 5 messages each, or terminate early by making a secret proposal at any time.} \\
\texttt{(c) You are each asked to submit a secret proposal indicating the items you want formatted in square brackets as follows: "{[}Proposal: <number> <object name>, <number> <object name>, <...>{]}"} \\
\texttt{(d) If your proposals are complementary, i.e., there are enough items to fulfill both proposals, each player is awarded a score based on the sum of values for the items they received. Otherwise, both of you get zero points.} \\
\texttt{(e) Your goal is to maximize the score you receive.} \\
\\ 
\texttt{Let us start.} \\
\\ 
\texttt{The set of available items is:} \\
\\ 
\texttt{1 comb, 2 whistles, 1 wrench, 1 statue} \\
\\ 
\texttt{Your secret value function is:} \\
\\ 
\texttt{comb: 0, whistle: 1, wrench: 0, statue: 8} \\
\\ 
\texttt{IMPORTANT: Your messages, unless it is the secret proposal, are directly transmitted to the other player, so do not include any response to the rules or text announcing your message. To make a secret proposal, use the indicated format. Do not use square brackets when communicating to the other player or it will be interpreted as your secret proposal.} \\
            }
        }
    }
    & \\ \\

    \theutterance \stepcounter{utterance}  
    & \multicolumn{4}{p{0.6\linewidth}}{
        \cellcolor[rgb]{0.9,0.9,0.9}{
            \makecell[{{p{\linewidth}}}]{
                \texttt{\tiny{[P1$\rangle$GM]}}
                \texttt{I take the statue and one whistle; you take the comb and the other whistle.} \\
            }
        }
    }
    & & \\ \\

    \theutterance \stepcounter{utterance}  
    & & & \multicolumn{2}{p{0.3\linewidth}}{
        \cellcolor[rgb]{0.9,0.9,0.9}{
            \makecell[{{p{\linewidth}}}]{
                \texttt{\tiny{[GM$|$GM]}}
                \texttt{continue} \\
            }
        }
    }
    & & \\ \\

    \theutterance \stepcounter{utterance}  
    & & & \multicolumn{2}{p{0.3\linewidth}}{
        \cellcolor[rgb]{0.9,0.9,0.9}{
            \makecell[{{p{\linewidth}}}]{
                \texttt{\tiny{[GM$|$GM]}}
                \texttt{I take the statue and one whistle; you take the comb and the other whistle.} \\
            }
        }
    }
    & & \\ \\

    \theutterance \stepcounter{utterance}  
    & & \multicolumn{4}{p{0.6\linewidth}}{
        \cellcolor[rgb]{0.9,0.9,0.9}{
            \makecell[{{p{\linewidth}}}]{
                \texttt{\tiny{[GM$\rangle$P2]}}
                \texttt{You have attempted this same task before and did not succeed. These are the plans you wrote after those attempts.} \\
\texttt{Trial \#1: New plan: I failed because I misinterpreted the game's termination condition. I assumed that once the other player submitted a secret proposal, I could simply submit my own proposal and the game would end with a score calculation. However, the rules state that a deal is only reached if the proposals are "complementary" (i.e., they partition the set of items without overlap). Since my proposal {[}1 statue, 1 wrench{]} did not account for the remaining items (the comb and whistles) or the other player's specific proposal, the game likely ended with zero points for both. My strategy was to propose a partial allocation rather than a complete, mutually exclusive partition. Next time, I must first determine the other player's likely proposal or value function, then construct a proposal that exactly covers the remaining items to ensure complementarity. I should aim for a proposal like "{[}1 statue, 1 wrench, 1 comb, 1 whistle{]}" if I intend to take everything, or a split that leaves exactly what the other player wants, ensuring no items are left unassigned or double{-}counted before submitting.} \\
\\ 
\texttt{Trial \#2: New plan: I failed because my secret proposal format was incorrect. The rules specify the format as "{[}Proposal: <number> <object name>, ...{]}", but I submitted "{[}Proposal: 1 statue, 1 whistle, 1 comb, 1 whistle{]}". This likely caused the parser to fail or misinterpret the list as a single item or a malformed string because I omitted the count for the second whistle (which is item \#2) and potentially the comma placement or the specific "number object" syntax for the second item. Next time, I will strictly follow the format "{[}Proposal: <count> <name>, <count> <name>...{]}" ensuring every item has a preceding number and proper comma separation, such as "{[}Proposal: 1 statue, 1 whistle, 1 comb, 1 whistle{]}" only if the parser expects "1 whistle" for the first and "1 whistle" for the second, but more likely I need to ensure the syntax matches exactly: "{[}Proposal: 1 statue, 1 whistle, 1 comb, 1 whistle{]}" might be valid if the parser handles duplicates, but the error suggests a structural issue. Looking closer at the error "rejected one of your messages because it violated the game rules", and my previous attempt in the prompt description (Trial \#1) mentioned a different proposal, but the actual failure here is the format. Wait, looking at the provided text: "You: {[}Proposal: 1 statue, 1 whistle, 1 comb, 1 whistle{]}". The rule says: "{[}Proposal: <number> <object name>, <number> <object name>, <...>{]}". My input has "1 whistle" twice. If the system expects unique identifiers or specific formatting, I might have failed. However, the most common format error in these tasks is missing the "number" prefix for the second item if the first was implicit, or simply the lack of a space after the colon. Let's re{-}read carefully: "{[}Proposal: <number> <object name>, <number> <object name>, <...>{]}". My input: "{[}Proposal: 1 statue, 1 whistle, 1 comb, 1 whistle{]}". This looks syntactically correct based on the template. Perhaps the issue is that I didn't account for the fact that there are *two* whistles and I listed them as two separate items, which is correct. Maybe the error was in the *previous* turn where I didn't use brackets? No, the prompt says "The other player has submitted a secret proposal. It is now your turn...". Ah, I see. In the "Your last attempt" section, the user's message was "I take the statue and one whistle; you take the comb and the other whistle." Then the GM said "The other player has submitted a secret proposal." Then the user said "{[}Proposal: 1 statue, 1 whistle, 1} \\
\\ 
\\ 
\texttt{You are playing a negotiation game in which you have to agree on how to divide a set of items among you and another player.} \\
\\ 
\texttt{Rules:} \\
\texttt{(a) You and the other player are given a set of items. Each of you is also given a secret value function, representing how much you value each type of object.} \\
\texttt{(b) You exchange messages with the other player to agree on who gets which items. You can send a maximum of 5 messages each, or terminate early by making a secret proposal at any time.} \\
\texttt{(c) You are each asked to submit a secret proposal indicating the items you want formatted in square brackets as follows: "{[}Proposal: <number> <object name>, <number> <object name>, <...>{]}"} \\
\texttt{(d) If your proposals are complementary, i.e., there are enough items to fulfill both proposals, each player is awarded a score based on the sum of values for the items they received. Otherwise, both of you get zero points.} \\
\texttt{(e) Your goal is to maximize the score you receive.} \\
\\ 
\texttt{Let us start.} \\
\\ 
\texttt{The set of available items is:} \\
\\ 
\texttt{1 comb, 2 whistles, 1 wrench, 1 statue} \\
\\ 
\texttt{Your secret value function is:} \\
\\ 
\texttt{comb: 1, whistle: 0, wrench: 4, statue: 5} \\
\\ 
\texttt{IMPORTANT: Your messages, unless it is the secret proposal, are directly transmitted to the other player, so do not include any response to the rules or text announcing your message. To make a secret proposal, use the indicated format. Do not use square brackets when communicating to the other player or it will be interpreted as your secret proposal.} \\
\\ 
\\ 
\texttt{I take the statue and one whistle; you take the comb and the other whistle.} \\
            }
        }
    }
    & \\ \\

    \theutterance \stepcounter{utterance}  
    & & & \multicolumn{4}{p{0.6\linewidth}}{
        \cellcolor[rgb]{0.9,0.9,0.9}{
            \makecell[{{p{\linewidth}}}]{
                \texttt{\tiny{[GM$\langle$P2]}}
                \texttt{{[}Proposal: 1 statue, 1 whistle, 1 comb, 1 whistle{]}} \\
            }
        }
    }
     \\ \\

    \theutterance \stepcounter{utterance}  
    & & & \multicolumn{2}{p{0.3\linewidth}}{
        \cellcolor[rgb]{0.9,0.9,0.9}{
            \makecell[{{p{\linewidth}}}]{
                \texttt{\tiny{[GM$|$GM]}}
                \texttt{proposal must include every item type only once} \\
            }
        }
    }
    & & \\ \\

    \theutterance \stepcounter{utterance}  
    & \multicolumn{4}{p{0.6\linewidth}}{
        \cellcolor[rgb]{0.95,0.95,0.95}{
            \makecell[{{p{\linewidth}}}]{
                \texttt{\tiny{[P2$\rightarrow$P1]}}
                \texttt{Post{-}episode Reflexion (not sent to GM):} \\
\texttt{New plan: The previous attempts failed because I either submitted a secret proposal too early without a counter{-}proposal to verify complementarity, or I included non{-}compliant text. My new strategy is to first send a direct, unbracketed verbal offer of a specific item split (e.g., "I take the statue and one whistle; you take the comb and the other whistle") to establish a potential agreement. I will wait for the opponent's response. If they accept or counter with a complementary offer, I will then submit the final secret proposal using the exact format `{[}Proposal: <number> <object name>, ...{]}` matching the agreed split. If they reject, I will adjust my verbal offer in the next message to be more attractive (e.g., offering the statue and both whistles) before attempting the final proposal again. Crucially, I will never use square brackets in any message until the final step, and I will ensure the verbal offer is a complete, direct statement without questions or explanations.} \\
            }
        }
    }
    & & \\ \\

    \theutterance \stepcounter{utterance}  
    & & & \multicolumn{4}{p{0.6\linewidth}}{
        \cellcolor[rgb]{0.95,0.95,0.95}{
            \makecell[{{p{\linewidth}}}]{
                \texttt{\tiny{[P1$\rightarrow$P2]}}
                \texttt{Post{-}episode Reflexion (not sent to GM):} \\
\texttt{New plan: I failed because I submitted a secret proposal while the game rules state that proposals are only submitted to finalize a deal, and the specific error context implies a formatting or timing issue with the proposal syntax relative to the game state. The rules specify the format as "{[}Proposal: <number> <object name>, <number> <object name>, <...>{]}". My previous attempt "{[}Proposal: 1 statue, 1 whistle, 1 comb, 1 whistle{]}" might have been rejected because the game expected a specific count for the duplicate item (whistle) or because I submitted a proposal when I should have been negotiating first, or simply because the parser was strict about the exact string matching including spaces or item identifiers. To ensure a legal response, I will first verify the exact item counts in the set (1 comb, 2 whistles, 1 wrench, 1 statue) and construct a proposal that explicitly lists each item with its correct quantity, ensuring the format is exactly "{[}Proposal: 1 comb, 2 whistles, 1 wrench, 1 statue{]}" if I am taking everything, or a valid subset if negotiating. Crucially, I will only submit a secret proposal if I have reached a consensus or if the game explicitly asks for it, rather than guessing the other player's move. I will start by communicating my intent to negotiate verbally to avoid premature termination or format errors.} \\
            }
        }
    }
     \\ \\

\end{supertabular}
}

\end{document}
